\chapter{Conclusiones y trabajo futuro}

El objetivo principal de este Trabajo de Fin de Grado ha consistido en diseñar e implementar una plataforma digital para la gestión colaborativa de explotaciones de olivar. El 
trabajo parte de la necesidad de organizar de forma estructurada la información asociada a parcelas, recintos, usuarios, entidades habilitadas y actividades agrícolas, en un 
contexto en el que la digitalización y la trazabilidad de la información agraria adquieren cada vez más importancia.

Como respuesta a esta problemática se ha desarrollado Olivaris, un prototipo funcional basado en una arquitectura cliente-servidor. El sistema incorpora un backend con API REST, 
autenticación mediante JWT, control de acceso por roles, gestión de usuarios y entidades, registro de parcelas y recintos, gestión de actividades agrícolas y soporte para información 
geoespacial. Además, se ha desarrollado una aplicación móvil que permite interactuar con las principales funcionalidades del sistema.

Los resultados obtenidos permiten concluir que la solución propuesta es técnicamente viable. Olivaris permite centralizar información relevante de una explotación agrícola, organizarla 
en torno a parcelas y actividades, y gestionar distintos perfiles de usuario con permisos diferenciados. De este modo, el sistema ofrece una base funcional para mejorar la organización 
de la información y facilitar la colaboración entre agricultores y entidades habilitadas.

Uno de los aspectos más relevantes del proyecto es la integración de distintos componentes tecnológicos en una solución coherente: backend, base de datos, información geoespacial, aplicación
móvil, mecanismos de seguridad, despliegue mediante contenedores y pruebas de rendimiento. Esta combinación permite que el sistema no se limite a una funcionalidad aislada, sino que constituya 
una base extensible sobre la que incorporar nuevas capacidades.

Asimismo, el diseño del sistema se ha planteado teniendo en cuenta el marco SIEX/REA/CUE, de forma que la estructura de la información pueda facilitar una posible adaptación futura a los requisitos 
de estas plataformas. Aunque no se ha implementado una integración efectiva con servicios oficiales, el trabajo realizado permite avanzar hacia una gestión agrícola más digital, estructurada y preparada 
para escenarios de interoperabilidad.

Las pruebas realizadas han permitido validar el comportamiento técnico del sistema en distintos escenarios de carga. Aunque esta evaluación no sustituye a una validación completa con usuarios reales, sí 
aporta una primera evidencia sobre el funcionamiento del backend y sobre la influencia que tiene la infraestructura de despliegue en el rendimiento de la aplicación.

En conjunto, Olivaris cumple los objetivos principales planteados al inicio del proyecto y demuestra que es posible construir una plataforma funcional para la gestión colaborativa del olivar utilizando 
tecnologías actuales, abiertas y extensibles.

\section{Limitaciones del sistema}

A pesar de los resultados obtenidos, el sistema presenta algunas limitaciones que deben tenerse en cuenta.

En primer lugar, la plataforma no incorpora todavía una integración efectiva con los servicios oficiales asociados al marco SIEX/REA/CUE. El diseño se ha orientado a facilitar una posible compatibilidad 
futura, pero la conexión real con dichas plataformas queda fuera del alcance implementado.

En segundo lugar, la aplicación no dispone de funcionamiento offline. Esta limitación es relevante en el contexto agrícola, ya que muchas parcelas pueden encontrarse en zonas con conectividad limitada. 
En el estado actual, el uso de la aplicación depende de la conexión con el backend.

Otra limitación es que la generación de informes se encuentra en una fase inicial. Aunque permite generarlos en base a las actividades realizadas, en una herramienta de este tipo sería útil poder generar 
informes en base a los criterios: parcela, campaña, producto aplicado o fecha.

Asimismo, el sistema se ha centrado principalmente en una aplicación móvil. Una versión web permitiría mejorar la consulta, revisión y administración de datos desde ordenador, especialmente para entidades 
habilitadas o usuarios que gestionen un número elevado de parcelas.

Finalmente, la evaluación se ha centrado en pruebas técnicas y de rendimiento. Sería necesario realizar pruebas con usuarios reales del sector agrícola para valorar la usabilidad de la aplicación, la claridad 
de los formularios y la adecuación de los flujos de trabajo a situaciones reales.

\section{Trabajo futuro}

A partir de las limitaciones identificadas, se plantean varias líneas de trabajo futuro.

Una de las más relevantes sería avanzar en la integración con el marco SIEX/REA/CUE, estudiando los formatos de intercambio de información, los requisitos técnicos y los mecanismos de autenticación necesarios 
para una futura conexión con servicios oficiales.

También resultaría importante incorporar funcionamiento offline y sincronización posterior, de forma que el usuario pueda registrar actividades en campo aunque no disponga de conexión en ese momento.

Otra línea de mejora sería ampliar la generación de informes, permitiendo filtrar y exportar información por parcela, campaña, producto aplicado, tipo de actividad o intervalo temporal. Esta funcionalidad aumentaría 
notablemente el valor práctico de la plataforma.

Además, sería conveniente desarrollar una interfaz web que complemente a la aplicación móvil. Esto permitiría ofrecer una experiencia más cómoda para tareas de revisión, administración, generación de informes y gestión 
de grandes volúmenes de información.

Por último, sería interesante realizar pruebas de usabilidad con agricultores, técnicos agrícolas o personal de cooperativas. Esta validación permitiría detectar problemas reales de uso, simplificar formularios y priorizar 
mejor las funcionalidades futuras.

\section{Conclusión final}

Olivaris ha permitido demostrar la viabilidad de una plataforma digital para la gestión colaborativa del olivar, integrando en una misma solución la gestión de usuarios, entidades, parcelas, recintos, actividades agrícolas, 
información geoespacial y control de permisos.

Aunque el sistema presenta limitaciones y todavía existen funcionalidades por desarrollar, el trabajo realizado constituye una base sólida y extensible. En este sentido, el TFG cumple su objetivo principal: construir un prototipo 
funcional y técnicamente coherente que contribuye a la digitalización de la gestión agrícola y que puede servir como punto de partida para futuras evoluciones hacia una herramienta más completa y cercana a las necesidades del sector.
